% ROZDZIAŁ 5 — plik włączany przez main.tex
\chapter{Struktura przestrzeni mierzalnych}
\label{ch:struktura-przestrzeni}

\section{Dlaczego sama definicja \texorpdfstring{$\sigma$}{sigma}-algebry nie wystarcza?}

Znamy już konstrukcję $\sigma$-algebr, miar i miary Lebesgue'a.
W dalszej teorii będziemy jednak przenosić miary przez odwzorowania,
budować przestrzenie produktowe i określać zbiory na podstawie
współrzędnych. Chcemy więc wiedzieć, jak z danych o prostych
zbiorach odzyskać mierzalność dowolnego odwzorowania.

W tym rozdziale przez $(X,\mathcal F)$ i $(Y,\mathcal G)$ rozumiemy
przestrzenie mierzalne. Nie zakładamy istnienia na nich miary,
jeśli nie podajemy jej wyraźnie. Zbiór $\mathbb N$ składa się z
dodatnich liczb całkowitych. Używamy oznaczeń
$\mathcal B(T)$ dla $\sigma$-algebry borelowskiej przestrzeni
topologicznej $T$ oraz $\mathcal F|_A=\{E\cap A:E\in\mathcal F\}$
dla śladu $\sigma$-algebry na $A\subseteq X$.

\section{Ślad \texorpdfstring{$\sigma$}{sigma}-algebry i podprzestrzenie}

Ślad $\mathcal F|_A$ zdefiniowaliśmy już w rozdziale~\ref{ch:przestrzenie-mierzalne}.
Teraz badamy jego zgodność z generowaniem $\sigma$-algebr
i topologią podprzestrzeni.


\begin{lemat}[Ślad i generowanie]
\label{lem:trace-generation}
Niech $X$ będzie zbiorem, $\mathcal E\subseteq\mathcal P(X)$,
$A\subseteq X$ i $\mathcal F=\sigma_X(\mathcal E)$.
Wtedy
\[
\mathcal F|_A=\sigma_A\bigl(\{A\cap E:E\in\mathcal E\}\bigr).
\]
\end{lemat}
\begin{proof}
Oznaczmy prawą stronę przez $\mathcal H$. Zbiory
$A\cap E$ dla $E\in\mathcal E$ należą do $\mathcal F|_A$;
ponieważ $\mathcal F|_A$ jest $\sigma$-algebrą na $A$,
mamy $\mathcal H\subseteq\mathcal F|_A$.
W drugą stronę rozważmy rodzinę
$\mathcal D=\{E\subseteq X:A\cap E\in\mathcal H\}$.
Zawiera ona $X$, gdyż $A\in\mathcal H$. Jeżeli
$E\in\mathcal D$, to
$A\cap(X\setminus E)=A\setminus(A\cap E)\in\mathcal H$.
Dla $E_n\in\mathcal D$ mamy
$A\cap\bigcup_nE_n=\bigcup_n(A\cap E_n)\in\mathcal H$.
Zatem $\mathcal D$ jest $\sigma$-algebrą na $X$, zawierającą
$\mathcal E$. Wynika stąd
$\mathcal F=\sigma_X(\mathcal E)\subseteq\mathcal D$, czyli
$\mathcal F|_A\subseteq\mathcal H$.
\end{proof}

\begin{twierdzenie}[Zbiory borelowskie podprzestrzeni]
\label{thm:borel-trace}
Niech $T$ będzie przestrzenią topologiczną i $A\subseteq T$
wyposażonym w topologię podprzestrzeni. Wtedy
$\mathcal B(A)=\mathcal B(T)|_A$.
\end{twierdzenie}
\begin{proof}
Z definicji topologii podprzestrzeni zbiory otwarte w $A$
mają postać $A\cap U$, gdzie $U$ jest otwarty w $T$.
Zastosowanie lematu \ref{lem:trace-generation} do rodziny
wszystkich zbiorów otwartych przestrzeni $T$ daje tezę.
\end{proof}

\begin{przyklad}[Ślad na zbiorze nieborelowskim]
Istnieje zbiór $A\subseteq\mathbb R$ nieborelowski
(wniosek \ref{cor:nonborel-lebesgue}). Mimo to możemy
rozpatrywać $(A,\mathcal B(\mathbb R)|_A)$.
To jest przestrzeń mierzalna; brak borelowskości $A$
w przestrzeni otaczającej nie uniemożliwia
zdefiniowania śladu. Nie wolno jednak utożsamiać
$\mathcal B(\mathbb R)|_A$ z rodziną
$\{E\in\mathcal B(\mathbb R):E\subseteq A\}$:
ta druga rodzina może nawet nie zawierać całego $A$.
\end{przyklad}

\section{Mierzalność odwzorowania i rodziny generujące}

Korzystamy z definicji mierzalności~\ref{def:measurable-map} i kryterium generatorów~\ref{thm:measurable-generator}, udowodnionych w rozdziale~\ref{ch:przestrzenie-mierzalne}. Ich podstawą są tożsamości dla przeciwobrazów: $f^{-1}(Y\setminus B)=X\setminus f^{-1}(B)$ oraz $f^{-1}(\bigcup_n B_n)=\bigcup_n f^{-1}(B_n)$. Nie wymagają one różnowartościowości ani surjektywności.

\begin{wniosek}[Prawostronny test mierzalności]
\label{cor:real-borel-tests}
Niech $(X,\mathcal F)$ będzie przestrzenią mierzalną,
a $f:X\to\mathbb R$ odwzorowaniem. Jest ono
$(\mathcal F,\mathcal B(\mathbb R))$-mierzalne wtedy i tylko
wtedy, gdy $\{x\in X:f(x)>q\}\in\mathcal F$
dla każdego $q\in\mathbb Q$.
\end{wniosek}
\begin{proof}
Rodzina prawych półprostych $(q,\infty)$, $q\in\mathbb Q$,
generuje $\mathcal B(\mathbb R)$: przez dopełnienia
otrzymujemy lewe półproste domknięte, a następnie
wszystkie otwarte przedziały. Stosujemy więc
kryterium generatorów~\ref{thm:measurable-generator}.
Warunki z lewymi półprostymi zostały już udowodnione
w twierdzeniu~\ref{thm:measurable-real}.
\end{proof}

\begin{twierdzenie}[Mierzalność ograniczenia]
\label{thm:restriction-meas-5}
Niech $(X,\mathcal F)$ i $(Y,\mathcal G)$ będą przestrzeniami
mierzalnymi, $f:X\to Y$ odwzorowaniem
$(\mathcal F,\mathcal G)$-mierzalnym oraz $A\subseteq X$.
Wówczas $f|_A:(A,\mathcal F|_A)\to(Y,\mathcal G)$ jest mierzalne.
\end{twierdzenie}
\begin{proof}
Dla każdego $B\in\mathcal G$ mamy
$(f|_A)^{-1}(B)=A\cap f^{-1}(B)\in\mathcal F|_A$.
Złożenia odwzorowań mierzalnych oraz mierzalność funkcji
ciągłych zostały już wykazane w twierdzeniu~\ref{thm:composition-measurable}.
\end{proof}

\section{Najmniejsza \texorpdfstring{$\sigma$}{sigma}-algebra, względem której funkcja jest mierzalna}

\begin{definicja}[$\sigma$-algebra generowana przez odwzorowanie]
\label{def:sigma-of-function}
Dla odwzorowania $f:X\to(Y,\mathcal G)$ określamy
\[
\sigma(f):=\{f^{-1}(B):B\in\mathcal G\}.
\]
Jeżeli $(f_i)_{i\in I}$ jest rodziną odwzorowań
$f_i:X\to(Y_i,\mathcal G_i)$, to
$\sigma(f_i:i\in I):=\sigma_X\bigl(\bigcup_{i\in I}\sigma(f_i)\bigr)$.
\end{definicja}

\begin{twierdzenie}[Minimalna struktura mierzalna]
\label{thm:initial-sigma}
Dla każdego $f:X\to(Y,\mathcal G)$ rodzina $\sigma(f)$
jest $\sigma$-algebrą na $X$ i jest najmniejszą
$\sigma$-algebrą, względem której $f$ jest mierzalne.
Dla rodziny $(f_i)_{i\in I}$ analogiczną własność
ma $\sigma(f_i:i\in I)$.
\end{twierdzenie}
\begin{proof}
Przeciwobraz całego $Y$ to $X$, przeciwobrazy zachowują
dopełnienia i przeliczalne sumy, więc $\sigma(f)$
jest $\sigma$-algebrą. Jeśli $f$ jest mierzalne
względem pewnej $\mathcal F$, to z definicji
$f^{-1}(B)\in\mathcal F$ dla każdego $B\in\mathcal G$,
toteż $\sigma(f)\subseteq\mathcal F$.
Dla rodziny odwzorowań powtarzamy argument dla
$\sigma$-algebry generowanej sumą wszystkich
rodzin przeciwobrazów.
\end{proof}

\begin{przyklad}[Rozróżnialność punktów]
Dla $X=\{a,b,c,d\}$ i funkcji $f:X\to\mathbb R$ o
$f(a)=f(b)=0$ oraz $f(c)=f(d)=1$ otrzymujemy
$\sigma(f)=\{\varnothing,\{a,b\},\{c,d\},X\}$.
Funkcja nie potrafi odróżnić $a$ od $b$ ani $c$ od $d$.
To wyjaśnia, dlaczego $\sigma(f)$ może być ściśle mniejsza
od $\mathcal P(X)$.
\end{przyklad}

\section{Iloczyn przestrzeni mierzalnych}

\begin{definicja}[$\sigma$-algebra produktowa]
\label{def:product-sigma}
Dla przestrzeni mierzalnych $(X,\mathcal F)$ i
$(Y,\mathcal G)$ definiujemy
\[
\mathcal F\otimes\mathcal G
:=\sigma_{X\times Y}\bigl(\{A\times B:A\in\mathcal F,\ B\in\mathcal G\}\bigr).
\]
Parę $(X\times Y,\mathcal F\otimes\mathcal G)$
nazywamy \emph{iloczynem przestrzeni mierzalnych}.
Nie wprowadzamy w tym miejscu miary produktowej;
jej konstrukcja należy do planowanego rozdziału 12.
\end{definicja}

\begin{twierdzenie}[Charakteryzacja przez współrzędne]
\label{thm:product-projections}
Niech $(X,\mathcal F)$ i $(Y,\mathcal G)$ będą
przestrzeniami mierzalnymi, a
$\pi_X:X\times Y\to X$, $\pi_Y:X\times Y\to Y$
rzutowaniami. Wówczas
$\mathcal F\otimes\mathcal G=\sigma(\pi_X,\pi_Y)$.
Ponadto odwzorowanie $h:Z\to X\times Y$ z przestrzeni
mierzalnej $(Z,\mathcal H)$ jest
$(\mathcal H,\mathcal F\otimes\mathcal G)$-mierzalne
wtedy i tylko wtedy, gdy oba odwzorowania
$\pi_X\circ h:Z\to X$ i $\pi_Y\circ h:Z\to Y$ są mierzalne.
\end{twierdzenie}
\begin{proof}
Ponieważ $\pi_X^{-1}(A)=A\times Y$ i
$\pi_Y^{-1}(B)=X\times B$, każda z rodzin
$\sigma(\pi_X)$, $\sigma(\pi_Y)$ zawiera się w
$\mathcal F\otimes\mathcal G$. Z drugiej strony
$A\times B=\pi_X^{-1}(A)\cap\pi_Y^{-1}(B)$,
zatem każdy generator $\mathcal F\otimes\mathcal G$
należy do $\sigma(\pi_X,\pi_Y)$. To dowodzi równości.
Drugą część otrzymujemy, sprawdzając przeciwobrazy
generatorów $A\times B$:
$h^{-1}(A\times B)=(\pi_X\circ h)^{-1}(A)
\cap(\pi_Y\circ h)^{-1}(B)$.
Konieczność wynika z mierzalności rzutowań i złożenia.
\end{proof}

\begin{twierdzenie}[Produkt borelowski przestrzeni drugoprzeliczalnych]
\label{thm:borel-product}
Niech $S,T$ będą przestrzeniami topologicznymi, z których
każda ma przeliczalną bazę topologii. Wtedy
\[
\mathcal B(S\times T)=\mathcal B(S)\otimes\mathcal B(T),
\]
gdzie na $S\times T$ rozważamy topologię produktową.
\end{twierdzenie}
\begin{proof}
Rzutowania $\pi_S,\pi_T$ są ciągłe, więc
$A\times T=\pi_S^{-1}(A)$ i $S\times B=\pi_T^{-1}(B)$
są borelowskie dla borelowskich $A,B$.
Stąd $\mathcal B(S)\otimes\mathcal B(T)\subseteq
\mathcal B(S\times T)$.
Wybierzmy przeliczalne bazy $\mathcal U,\mathcal V$.
Wszystkie $U\times V$, gdzie $U\in\mathcal U$ i
$V\in\mathcal V$, należą do
$\mathcal B(S)\otimes\mathcal B(T)$ i tworzą
przeliczalną bazę produktu. Każdy zbiór otwarty
w $S\times T$ jest więc przeliczalną sumą takich
prostokątów, czyli należy do $\mathcal B(S)\otimes\mathcal B(T)$.
Generując $\sigma$-algebrę, otrzymujemy zawieranie odwrotne.
\end{proof}

\begin{wniosek}[Borelowskie odwzorowania wektorowe]
\label{cor:coordinate-measurable}
Niech $(X,\mathcal F)$ będzie przestrzenią mierzalną,
$f=(f_1,\ldots,f_d):X\to\mathbb R^d$ i $d\ge1$.
Odwzorowanie $f$ jest $(\mathcal F,\mathcal B(\mathbb R^d))$-mierzalne
wtedy i tylko wtedy, gdy każda $f_j:X\to\mathbb R$
jest $(\mathcal F,\mathcal B(\mathbb R))$-mierzalna.
\end{wniosek}
\begin{proof}
Z twierdzenia \ref{thm:borel-product}, zastosowanego
indukcyjnie, wynika
$\mathcal B(\mathbb R^d)=\bigotimes_{j=1}^d\mathcal B(\mathbb R)$.
Wykorzystujemy teraz twierdzenie \ref{thm:product-projections}.
\end{proof}

\begin{uwaga}[Dlaczego potrzebna jest przeliczalna baza?]
Dla dowolnych przestrzeni topologicznych $S,T$ zachodzi
$\mathcal B(S)\otimes\mathcal B(T)\subseteq\mathcal B(S\times T)$,
ponieważ rzutowania są ciągłe. Równość może jednak zawieść.
Dla przykładu niech $D$ będzie zbiorem o mocy większej niż
$\mathfrak c=|\mathbb R|$, wyposażonym w topologię dyskretną.
Wtedy przekątna $\Delta=\{(x,x):x\in D\}$ jest zamknięta
w $D\times D$, ale nie należy do
$\mathcal P(D)\otimes\mathcal P(D)$.
Istotnie, każdy element $\sigma$-algebry generowanej
przez dowolną rodzinę zbiorów należy do $\sigma$-algebry
generowanej przez pewną jej przeliczalną podrodzinę:
rodzina zbiorów o tej własności sama jest $\sigma$-algebrą.
Gdyby $\Delta$ była mierzalna produktowo, istniałyby
prostokąty $A_n\times B_n$, $n\ge1$, generujące
$\sigma$-algebrę zawierającą $\Delta$.
Dla różnych $x,y\in D$ punkty $(x,x)$ i $(x,y)$
muszą być rozróżnione przez któryś z tych prostokątów,
inaczej wszystkie zbiory z wygenerowanej $\sigma$-algebry
zawierałyby oba albo żaden z nich. Ponieważ pierwsza
współrzędna obu punktów jest taka sama, pewien $B_n$
musi zawierać dokładnie jeden z punktów $x,y$.
Odwzorowanie $x\mapsto(\chi_{B_n}(x))_{n\ge1}$
byłoby więc iniekcją $D\to\{0,1\}^{\mathbb N}$, co
przeczy $|D|>\mathfrak c$.
\end{uwaga}

\section{Przestrzenie standardowe borelowskie}

\begin{definicja}[Przestrzeń polska i standardowa borelowska]
\label{def:standard-borel}
Przestrzeń topologiczną nazywamy \emph{polską}, jeśli
jej topologię wyznacza pewna metryka zupełna i ośrodek
przeliczalny, czyli istnieje w niej przeliczalny zbiór
gęsty. Przestrzeń mierzalną $(X,\mathcal F)$ nazywamy
\emph{standardową borelowską}, jeżeli istnieje przestrzeń
polska $P$ i bijekcja $b:X\to P$, dla której
$b$ oraz $b^{-1}$ są mierzalne względem
$\mathcal F$ i $\mathcal B(P)$. Taką bijekcję
nazywamy \emph{izomorfizmem przestrzeni mierzalnych}.
\end{definicja}

\begin{uwaga}[Uwaga o definicji]
Równoważnie można zażądać, aby $\mathcal F$ była
$\sigma$-algebrą borelowską pewnej polskiej topologii na $X$.
Istotny jest charakter \emph{mierzalny}, a nie zgodność
z jakąkolwiek wcześniej wybraną topologią. W definicji
żądaliśmy jawnie mierzalności odwrotności; dla ogólnych
przestrzeni mierzalnych sama mierzalna bijekcja nie wystarcza.
\end{uwaga}

\begin{przyklad}[Klasyczne przestrzenie standardowe]
Przestrzenie $(\mathbb R^d,\mathcal B(\mathbb R^d))$
($d\ge1$), $(\mathbb N,\mathcal P(\mathbb N))$ oraz
$(\{0,1\}^{\mathbb N},\mathcal B(\{0,1\}^{\mathbb N}))$
są standardowe borelowskie. W pierwszych dwóch przypadkach
korzystamy ze zwykłej metryki euklidesowej lub metryki
dyskretnej. W trzecim stosujemy metrykę
$\rho(u,v)=\sum_{n\ge1}2^{-n}|u_n-v_n|$ z rozdziału
\ref{ch:lebesgue}. Zupełność wynika stąd, że w ciągu
Cauchy'ego każda współrzędna musi od pewnego miejsca
być stała; przeliczalnym zbiorem gęstym są ciągi
ostatecznie równe zeru. Metryka $\rho$ wyznacza topologię
produktową, więc mamy przestrzeń polską.
\end{przyklad}

\begin{twierdzenie}[Przeliczalne generowanie standardowych przestrzeni]
\label{thm:standard-countably-generated}
Jeżeli $(X,\mathcal F)$ jest standardową przestrzenią
borelowską, to istnieje przeliczalna rodzina
$\mathcal E\subseteq\mathcal F$ taka, że
$\mathcal F=\sigma_X(\mathcal E)$. Ponadto dla
każdej pary różnych punktów $x,y\in X$ istnieje
$E\in\mathcal E$ zawierający dokładnie jeden z nich.
\end{twierdzenie}
\begin{proof}
Niech $b:X\to P$ będzie izomorfizmem na przestrzeń polską.
Każda przestrzeń metryczna o przeliczalnym zbiorze gęstym
ma przeliczalną bazę: bierzemy kule o środkach w tym zbiorze
i dodatnich promieniach wymiernych. Niech
$\mathcal U=\{U_n:n\ge1\}$ będzie taką bazą $P$.
Rodzina $\mathcal E=\{b^{-1}(U_n):n\ge1\}$
jest przeliczalna i generuje $\mathcal F$, gdyż
$\mathcal B(P)=\sigma_P(\mathcal U)$, a przeciwobrazy
zachowują dopełnienia i sumy przeliczalne.
Jeśli $x\ne y$, to $b(x)\ne b(y)$; w przestrzeni
metrycznej istnieje zbiór otwarty zawierający $b(x)$,
lecz nie $b(y)$. Z własności bazy można wybrać
$U_n$ o tej samej własności, a $b^{-1}(U_n)$
oddziela $x$ od $y$.
\end{proof}

\begin{przyklad}[Uzupełnienie miary zmienia typ przestrzeni]
\label{ex:completion-not-standard}
Przestrzeń $(\mathbb R,\mathcal B(\mathbb R))$ jest
standardowa borelowska, ale
$(\mathbb R,\mathcal L_1)$ nie jest standardowa borelowska.
Z twierdzenia \ref{thm:standard-countably-generated}
każda standardowa $\sigma$-algebra jest przeliczalnie
generowana. Każda $\sigma$-algebra generowana przez
przeliczalną rodzinę ma co najwyżej $\mathfrak c$
elementów (kodowanie zbiorów borelowskich drzewami,
jak w dowodzie wniosku \ref{cor:nonborel-lebesgue}).
Tymczasem $\mathcal L_1$ zawiera wszystkie podzbiory
zbioru Cantora, a więc co najmniej $2^{\mathfrak c}>
\mathfrak c$ elementów. Nie może być przeliczalnie
generowana. Przykład pokazuje, że \emph{zupełność miary}
i \emph{standardowość borelowska} są innymi własnościami.
\end{przyklad}

\section{Struktura mierzalna a miara obrazowa}

Przypomnijmy, że jeśli $\mu$ jest miarą na $(X,\mathcal F)$,
a $f:(X,\mathcal F)\to(Y,\mathcal G)$ jest mierzalne,
to \emph{miara obrazowa} $f_\#\mu$ określona jest wzorem
$(f_\#\mu)(B)=\mu(f^{-1}(B))$, $B\in\mathcal G$
(definicja z rozdziału \ref{ch:miary-wlasnosci}).
To poprawne określenie miary, ponieważ przeciwobrazy
zachowują rozłączność i przeliczalne sumy.

\begin{twierdzenie}[Funktorialność miar obrazowych]
\label{thm:pushforward-composition}
Niech $\mu$ będzie miarą na $(X,\mathcal F)$,
$f:(X,\mathcal F)\to(Y,\mathcal G)$ i
$g:(Y,\mathcal G)\to(Z,\mathcal H)$ niech będą mierzalne.
Wtedy
\[
(g\circ f)_\#\mu=g_\#(f_\#\mu).
\]
\end{twierdzenie}
\begin{proof}
Dla każdego $C\in\mathcal H$ z definicji i
rachunku przeciwobrazów
$(g\circ f)_\#\mu(C)=\mu(f^{-1}(g^{-1}(C)))
=(f_\#\mu)(g^{-1}(C))=g_\#(f_\#\mu)(C)$.
Równość miar oznacza tu równość ich wartości
na wszystkich zbiorach mierzalnych.
\end{proof}

\begin{przyklad}[Miara obrazowa przez funkcję stałą]
Niech $f:X\to Y$ będzie funkcją stałą o wartości
$y_0\in Y$ i niech $\mu$ będzie miarą na $(X,\mathcal F)$.
Wtedy $(f_\#\mu)(B)=\mu(X)$, gdy $y_0\in B$, oraz
$(f_\#\mu)(B)=0$, gdy $y_0\notin B$. Jeśli $\mu(X)<\infty$, można zapisać
$f_\#\mu=\mu(X)\delta_{y_0}$, gdzie $\delta_{y_0}$
wyraża miarę Diraca w punkcie $y_0$.
Gdy $\mu(X)=\infty$, pozostajemy przy podanym wzorze
zbiorowym: wyrażenie $\infty\cdot0$ nie jest w zwykłej
arytmetyce rozszerzonej określone.
\end{przyklad}

\section{Podsumowanie i powiązania}

Rodziny generujące upraszczają testy mierzalności.
Ślad opisuje przestrzenie indukowane, a produkt
sprowadza mierzalność wektorową do współrzędnych.
Ważnym rozróżnieniem jest różnica między przestrzenią
standardową borelowską a uzupełnioną przestrzenią
z miarą Lebesgue'a.

W następnym rozdziale wykorzystamy te narzędzia do
mierzalności funkcji rzeczywistych, ich aproksymacji
funkcjami prostymi i badania granic ciągów funkcji.